\section{Propagated campaign}
\label{sec:campaign}

Sections~\ref{sec:benchmark} and~\ref{sec:controller} are first-order
statements about a fixed reference trajectory. This section propagates.

\subsection{Screening}

Each augmented variational solve costs about one propagation, and a propagated
comparison costs several, so candidates were screened before they were flown.
The screening instrument is the forced variational solution, which in the
previous campaign reproduced the sign of the propagated ratio on \ph{n} of
\ph{n} resolved comparisons of a \ph{n}-orbit panel. All candidates share one
reference trajectory and one gradient inside a single augmented integration and
differ only through their $\Delta\mathbf a$, so a panel of \ph{n} orbits is
solved once per orbit with every candidate as a separate forcing channel rather
than once per candidate. The selection rule, the panel and the number of
candidates allowed through --- two --- were fixed before the screen ran.
\ph{list the candidates and the outcome.} \wpref{WP13}

\paragraph{One caution about the screen, stated because it is load-bearing.}
The previous campaign's calibration of this instrument was established on
policies that were \emph{not} constructed by minimizing $J$. \Asign\ is, and so
are the controllers distilled from it. Scoring a candidate with the same
functional it was optimized against ranks candidates fairly but does not test
whether $J$ predicts $E$ for that class of schedule --- it assumes it. That
assumption is what the sign check below tests, which is why a failure there is
treated as a statement about the instrument rather than as noise, and why the
screening escape of Section~\ref{sec:setup-prereg} bounds how long it may be
diagnosed.

\subsection{What is propagated}

Five policies are propagated at each campaign point, and two of them are there
for reasons worth stating. \Asign\ is propagated because the first-order
benchmark has to be tested as a trajectory statement, not assumed to be one,
and because it is the denominator of the capture fraction below. \Rint\ is
propagated because it is the strongest deployable schedule available before
this paper, so beating \Fop\ while losing to \Rint\ would not be a result.

\begin{center}\small
\begin{tabular}{@{}ll@{}}
  \toprule
  Propagated & \Cplan, \Cfb\ (or the two screened candidates), \Asign, \Fop, \Rint \\
  Not propagated & \Fenv, constructed as a lower envelope over the constant
    family sweep \\
  \bottomrule
\end{tabular}
\end{center}

Each of the four calibrated policies --- the two screened controllers,
\Asign\ and \Rint\ --- is brought onto the anchor $B_{\mathrm{tot}}$ by the
iteration of Section~\ref{sec:setup-budgets}, which is why the campaign
propagates more arcs than it compares: \ph{value} propagations per policy per
orbit, \ph{n} in total. \Fop\ is the anchor and is not calibrated.
Table~\ref{tab:beta1} reports the achieved match alongside each verdict, and no
comparison is read where the match falls outside the declared band.

Arcs from the previous campaign are reused only where their realized work can
be matched under the same convention; where it cannot, the arc is propagated
again. Every reused arc is flagged as such in the record, because the earlier
campaign matched a different budget definition and a silent reuse would charge
that difference to the method. \ph{n} arcs were reused and \ph{n} were
recomputed.

\subsection{Equal realized work at $\beta=1$}

\begin{table}[!htbp]
  \centering\small
  \caption{The controller against both constant-degree comparators at
  $\beta=1$, matched on realized total quadratic work including the
  controller's own overhead $B_+$. Counts are per-orbit comparisons clearing
  the summed reference-inclusive envelope; $\rho$ is the median of per-orbit
  ratios over all orbits of the design, not over the resolved subset. Every
  candidate sits on the anchor $B_{\mathrm{tot}}$ to within \ph{pct}; the
  achieved match is archived per orbit and its worst case quoted here.
  \ph{fill from WP14}}
  \label{tab:beta1}
  \begin{tabular}{@{}lcccccc@{}}
    \toprule
    & \multicolumn{3}{c}{Design A} & \multicolumn{3}{c}{Design B} \\
    \cmidrule(lr){2-4}\cmidrule(lr){5-7}
    Comparator & wins & losses & $\rho$ & wins & losses & $\rho$ \\
    \midrule
    \Fop\ (deployable)   & \ph{n} & \ph{n} & \ph{value} & \ph{n} & \ph{n} & \ph{value} \\
    \Rint\ (deployable)  & \ph{n} & \ph{n} & \ph{value} & \ph{n} & \ph{n} & \ph{value} \\
    \Fenv\ (post hoc)    & \ph{n} & \ph{n} & \ph{value} & \ph{n} & \ph{n} & \ph{value} \\
    \bottomrule
  \end{tabular}
\end{table}

\ph{two paragraphs} The verdict against \Fop\ and \Rint\ is the claim; the row
against \Fenv\ is reported because a reader is entitled to know how much of the
constant family's potential the deployable comparator leaves unused, and it is
not a policy anyone could select in advance.

\begin{figure}[!tbp]
  \centering
  \ph{figures/fig\_oa\_beta1\_per\_orbit.pdf}
  \caption{Orbit-by-orbit ratio at $\beta=1$, sorted within each design.
  Colour records each orbit's own resolution verdict, so grey bars are
  comparisons the experiment cannot decide rather than ties.}
  \label{fig:beta1}
\end{figure}

\subsection{Capture fraction}

The controller is measured against the benchmark it approximates:

\begin{equation}
  f \;=\; \frac{E_{\Fop}-E_{\Cplan}}{E_{\Fop}-E_{\Asign}} ,
  \qquad\text{defined only where}\quad
  E_{\Asign} < E_{\Fop}
  \;\text{ and }\;
  M_{\mathrm{res}}(\Fop,\Asign) > 1 .
  \label{eq:capture}
\end{equation}

Three details are fixed in advance. The benchmark \Asign\ is not deployable and
carries no overhead, so it receives the full $B_{\mathrm{tot}}$ while \Cplan\ receives $B_{\mathrm{tot}}-B_+$; that asymmetry is deliberate, and it is why
$f$ measures the price of deployability rather than only the quality of the
approximation. Both terms are propagated errors $E$, never predictions, so $f$
does not exist before this section. And $f>1$ is possible without being an
error: \Asign\ is a first-order optimum on the reference arc and can be beaten
on the arc it is actually flown along, which is what
Section~\ref{sec:campaign-fixedpoint} tests separately.

The validity condition is not cosmetic. Where the benchmark does not beat the constant
degree the denominator is negative and $f$ has no interpretation; where the two
are separated by less than the numerical envelope it is a ratio of two
unresolved quantities and diverges. Orbits failing either test are reported as
\emph{not applicable} and counted, never coerced to zero or to a large number.
\ph{n} of \ph{n} orbits are eligible. \ph{one paragraph with the median $f$ by
design and budget.}

\begin{figure}[!tbp]
  \centering
  \ph{figures/fig\_oa\_capture.pdf}
  \caption{Capture fraction across the budget grid, both designs.}
  \label{fig:capture}
\end{figure}

\subsection{Across the budget grid}

\ph{one paragraph plus table} $\beta\in\{0.50,0.75,1.00,1.25,1.50\}$. The
compute-starved end is where every deployable alternative lost in the previous
campaign, so it is reported first and separately. \ph{state the verdict at
$\beta=0.50$ explicitly, whichever way it goes.}

\dnote{H4 and H5 are decided here and are separate hypotheses on purpose: H5
($\beta=0.5$) can be rejected without touching H4. Report a rejection as a
rejection.}

\subsection{Across populations}

\ph{one paragraph plus table} Designs A, B, C; the five geometry strata; the
wide-elliptic population, which is the regime where the previous paper's radial
endpoint \emph{won}, so a trajectory-aware controller that fails there is
regime-bound; and the low-perilune subset, where the linearization behind both
the benchmark and the planner is weakest and only signs are read.

\subsection{Cost ladder}

\begin{table}[!htbp]
  \centering\small
  \caption{What ``equal budget'' equals, down the four costs of
  Section~\ref{sec:setup-budgets}, as controller over comparator. Only the
  primary row is matched by construction. \ph{fill from WP16}}
  \label{tab:cost-ladder}
  \begin{tabular}{@{}lccl@{}}
    \toprule
    Cost & A & B & \\
    \midrule
    Realized total quadratic work $B_2$ & \ph{value} & \ph{value} & matched by construction \\
    \quad of which controller overhead $B_+$ & \ph{pct} & \ph{pct} & pilot arc, probe, planning \\
    Nominal per-call work $B_1$ & \ph{value} & \ph{value} & reported for comparability only \\
    Right-hand-side calls made & \ph{value} & \ph{value} & \\
    End-to-end wall time & \multicolumn{2}{c}{\ph{value}} & on \ph{n} timed orbits \\
    Measured kernel time $B_3$ & \multicolumn{2}{c}{\ph{value}} & idle machine, three repeats \\
    \bottomrule
  \end{tabular}
\end{table}

\ph{one paragraph} The overhead row is the one to read first: a controller that
wins on error while spending its margin on its own planning has not won.

\subsection{Does the first-order benchmark survive propagation?}
\label{sec:campaign-fixedpoint}

\Asign\ optimizes a linearization about a fixed reference arc. When its
schedule is flown, the trajectory changes, $\Phi$ changes, and the local defect
changes with them, so the schedule that was optimal for the reference arc need
not be optimal for its own. The fixed-point test propagates the schedule,
re-tabulates along the resulting arc, re-solves, and propagates again.
\ph{one paragraph: how far the schedule moves, how far the error moves, and
after how many iterations.} \wpref{WP17}

\dnote{This is the test that decides whether \Asign\ is a
\emph{trajectory-level} benchmark or a \emph{reference-arc} one. If the gain
does not survive, say so in these words and rename the quantity throughout;
that outcome is a strong negative result about signed allocation and is worth
stating cleanly rather than burying.}

\subsection{Controls}

Each was drawn under a rule fixed before the runs it selects, and each targets
a way the result could be an artefact rather than a measurement.

\begin{table}[!htbp]
  \centering\small
  \caption{Registered controls. \ph{fill}}
  \label{tab:controls}
  \begin{tabular}{@{}p{0.24\textwidth}p{0.30\textwidth}p{0.38\textwidth}@{}}
    \toprule
    Control & What it tests & Outcome \\
    \midrule
    Reference degree $300\to600$ & whether large ratios are
      reference-adjacency artefacts rather than measurements & \ph{outcome} \\
    Decision-grid parity & that \Asign\ and \Cplan\ switch on the same grid, so
      the benchmark's margin is information and not resolution & \ph{outcome} \\
    Degree-cap audit & orbits whose schedule reaches the reference degree,
      where its force model \emph{is} the reference model & \ph{outcome} \\
    Phase-shifted initial states & whether the signed allocation is fitted to
      one arc's cancellation pattern & \ph{outcome} \\
    Time-indexed against altitude-binned schedule & whether the gain belongs to
      the schedule's form & \ph{outcome} \\
    Decision-grid coarsening & whether it belongs to the decision cadence &
      \ph{outcome} \\
    Information-source audit & that no \Cplan\ run read the reference field or
      the reference arc & \ph{outcome} \\
    \bottomrule
  \end{tabular}
\end{table}
